\section{A fully specified polynomial with 854 variables}
\label{sec:explicit-circuit}

This section gives a finite formula for an actual polynomial, not just
an existence argument. Its degree is much higher than four. Arithmetic
expressions in this section introduce no quantified coordinates unless
explicitly listed. The accompanying JSON circuit makes every addition,
multiplication, power, residual numerator, and nonzero condition literal.

\subsection{The finite list of rational equations}

Use the curve and constants of Section~\ref{sec:source-structure}, and
write $\sigma$ for the automorphism of $K=\Q(\sqrt2)$ sending
$\sqrt2$ to $-\sqrt2$. For a nonidentity point $Q=(x,y)$ define
\[
 f(Q)=1/x,\qquad z(Q)=-x/y.
\]
All the denominators occurring here and below are required to be
nonzero. A field inverse is written using its rational norm, so both
rational coefficient formulas are specified.

There are six root points $A,S,T,Z,J_0,J$, and five private points
$Q_{\ell,0},\ldots,Q_{\ell,4}$ for each $1\leq\ell\leq41$.
Their total number is $211$. For each point use the three rational
coordinates $(a_0,a_1,C_Q)$ and the common chart parameter
$j\in\{0,\ldots,211\}$ of
Proposition~\ref{prop:quadratic-square-root}. Namely put
$a=a_0+a_1\sqrt2$, write $a^3-Da=A_Q+B_Q\sqrt2$, reconstruct
$b=b_0+b_1\sqrt2$, impose
\[
 C_Q^2 V_Q-(B_Q/2+jV_Q)^2=0,
 \qquad V_Q=\frac{C_Q^2-A_Q-jB_Q}{j^2-2},
\]
and use the double $Q=[2](a,b)$ in every subsequent expression.
Require $C_Q\ne0$, $N_{K/\Q}(b)\ne0$, and nonzero norms for
the abscissa and ordinate of the output point. The source-chart
equivalence and its inverse maps are proved in
Section~\ref{sec:finite-fibre-application}.

Set $\beta=X^2+2$. Introduce rational proxies $s,\tau,z$ and impose
\begin{equation}\label{eq:circuit-lucas}
 \tau^2-s(s-2)\tau z+(s-2)^2z^2-1=0.
\end{equation}
Use one rational $\theta$, with $\varepsilon=\theta^2-2$.
The letter $z$ without an argument in \eqref{eq:circuit-lucas}
is a proxy; $z(Q)$ is the specified elliptic rational function.

The finite reciprocal logarithms are
\[
 \mathcal L_L(Q)=z(Q)G_{L/4-1}(Df(Q)^2),
 \qquad L\in\{16,64\},
\]
where $G_J(W)=1+\sum_{r=1}^J a_rW^r$,
$a_1=-2/5$, and $a_r=a_{r-1}(4r-5)/(4r+1)$.
Their nonvanishing on the required rational point charts is proved
in the finite-fiber application.

We next specify the norm operation used in every contact. Given a
nonzero rational $r$, an ordered list of tests $t_i$ with positive
orders $L_i$, and $L_1=1$, form the normalized leaves
$1,t_1/r,t_2/r^{L_2},\ldots$. First pair $1,t_1/r$ by
$N_r(a,b)=a^2+rb^2$. Group all remaining tests by increasing order,
preserving their original order within each group; pair successive
normalized leaves and carry an unpaired last leaf. At each higher
level pair successive nodes by $N_r$ and carry any last node until
one root $W_r$ remains. The input of the pole equations is $b=rW_r^2$.
This prescription is finite and fixes the order of every operation.

For a rational input $b$, a pole occurrence has five coordinates
$t,e_p,p,q,h_p$. Define
\[
 U=\frac{b(1-bt^2)}{1+bt^2},\qquad B=17b,
 \qquad H=e_p(B-U^2),
\]
and impose
\begin{equation}\label{eq:circuit-pole}
 p^2=B(1-e_p)(1+3e_p),\qquad
 Bq^2=(h_p^2-BH)(h_p^2+3BH).
\end{equation}
The denominator $1+bt^2$ cannot vanish on the original pole equations;
its explicit nonzero condition is retained here.

For the height occurrence take
\[
 r_0=\beta(\beta-1)(\beta-2401)(32\beta+343)(2\beta-1)
\]
and the six rational coefficients of
$f(J_0),f(S),z(S)/z(J_0)$, in that order. Their orders are
$(1,1,1,1,48,48)$. Replace the resulting pole input by $c r_0W_{r_0}^2$,
where the fixed finite choice is
\[
 c\in\mathcal C=\left\{\pm\prod_{p\in I}p:
 I\subseteq\{2,3,5,7,11,13,17\}\right\}.
\]
There are $256$ such integers, covering the source's eight Boolean
choices exactly.

Introduce a nonzero coordinate $e$, and let $a_1,\ldots,a_9$ be
$e$ followed by the eight rational coefficients of
$ex(A),ex(S),ex(T),ex(Z)$. Introduce two contact coordinates $u_c,v_c$.
With $W_0=25878772920$, the forty-one contacts are
\[
 \frac{v_c}{32W_0},\quad u_c,\quad u_c-v_c,\quad
 u_c-2401v_c,\quad u_c+\frac{343}{32}v_c,
\]
and, for each $i=1,\ldots,9$, on putting $A_i=47027a_i+3$,
\[
 u_c-A_iv_c,\quad u_c-A_iv_c+v_c,\quad
 u_c-A_iv_c+2401v_c,\quad
 u_c-A_iv_c-\frac{343}{32}v_c.
\]
Write these contacts as $r_1,\ldots,r_{41}$ in the displayed order.

For contact $\ell$, the five target points are
$(4P_*,4A,4S,4T,4Z)$. Subtract them respectively from the five
private points to obtain $Y_{\ell,i}$, using the affine subtraction
formula and its nonzero abscissa-difference denominator. The ordered
twenty field tests are the following:
\begin{enumerate}
\item $f(Q_{\ell,i})$ for $0\leq i\leq4$;
\item $f(J)^\sigma$;
\item $f(Y_{\ell,i})^\sigma$ for $0\leq i\leq4$;
\item $\mathcal L_{16}(Q_{\ell,1})/\mathcal L_{64}(Q_{\ell,0})-\beta$,
and the analogous three ratios with numerator precision $64$ and
proxies $s,\tau,z$;
\item $z(Y_{\ell,i})^\sigma/z(J)^\sigma$ for $0\leq i\leq4$.
\end{enumerate}
Replace each field test $A+B\sqrt2$ by the two rational tests
$\varepsilon B,A+\theta B$, and prepend $\varepsilon$. The
resulting forty-one orders are
\[
 (\underbrace{1,\ldots,1}_{23},16,16,
 \underbrace{64,\ldots,64}_{8},16,16,
 \underbrace{64,\ldots,64}_{6}).
\]
Apply the norm prescription at $r_\ell$ and then one copy of
\eqref{eq:circuit-pole}. This makes $42$ pole occurrences in total.

Finally, write $f(S)=f_0+f_1\sqrt2$, set
$F=f_0^2+f_0f_1-f_1^2$, and let $C=3^{2^{24}}$ be the fixed
anchor constant. The three-witness anchor is
\[
 u_a^2+6v_a^2-2w_a^2-3F^2+2C^4=0.
\]
We next absorb its private coordinate $u_a$ and all denominator
conditions into one coordinate, without increasing their combined count.

\subsection{Eliminating the denominator witness}

For fixed $(j,c)$, clear the denominators of all rational equations
other than the anchor, and call their polynomial numerators
$E_1,\ldots,E_{296}$. There are $211$ point equations, $84$ pole
equations, and one Lucas equation. During the construction record every
nonzero numerator required for a division, including rational norms
for divisions in $K$, and the explicit point and coordinate guards.
Let $G$ be their product, with repeated factors omitted if desired.
All the rational functions used are defined whenever $G\ne0$.

Write
\[
 6v_a^2-2w_a^2-3F^2+2C^4=N/B
\]
with integer polynomial numerator $N$ and denominator $B$. Choose
$F=R_0/S_0$ in the rational-expression construction. Then
\[
 B=S_0^2,\qquad
 N=S_0^2(6v_a^2-2w_a^2+2C^4)-3R_0^2.
\]
The actual circuit's cleared anchor has exactly this form.
All denominator factors are included among the recorded conditions.
The three anchor coordinates occur nowhere in $G$ or the other
residuals, and $B$ is independent of them.

The anchor always has a solution with $u_a\ne0$ whenever it has any
solution. Indeed its three witnesses cannot all vanish, since
$3F^2=2C^4$ is impossible for rational $F/C^2$. If $u_a=0$ and
$w_a\ne0$, replace $(u_a,w_a)$ by $(4w_a,3w_a)$; this preserves
$u_a^2-2w_a^2$. If $u_a=w_a=0$ and $v_a\ne0$, replace
$(u_a,v_a)$ by $(-12v_a/7,-5v_a/7)$; this preserves
$u_a^2+6v_a^2$. In either case $u_a$ becomes nonzero without
changing any outside equation.

Rescale $t_a=S_0u_a$, so that the cleared anchor is $t_a^2+N=0$.
Replace its private coordinate by $U$ and put
\[
 \Psi=1+NG^2U^2.
\]
Define the branch polynomial
\begin{equation}\label{eq:854-branch}
 F_{j,c}=\sum_{i=1}^{296}E_i^2+\Psi^2,
\end{equation}
and the required integer polynomial
\begin{equation}\label{eq:854-polynomial}
 P_{854}=\prod_{j=0}^{211}\ \prod_{c\in\mathcal C}F_{j,c}.
\end{equation}

\begin{proposition}\label{prop:854-explicit}
The polynomial \eqref{eq:854-polynomial} uses 854 total variables
and has the same rational input projection as $\Src$. A supplied
integer arithmetic circuit specifies it completely. Its degree is
at most $63570781474816$, and its coefficient height is at most
$2^{757866260454400}$.
\end{proposition}
\begin{proof}
If $GU=0$, then $\Psi=1$, so \eqref{eq:854-branch} has no rational
zero. Otherwise $G\ne0$ also implies $S_0\ne0$, and
\[
 \Psi=(GU)^2\bigl((1/(GU))^2+N\bigr).
\]
A zero therefore gives an original anchor witness
$u_a=1/(GS_0U)$. Conversely, use the preceding transformations to
choose $u_a\ne0$, then put $U=1/(GS_0u_a)$. This preserves the
existential input projection. It need not preserve a fiber in which
both of the other private anchor witnesses have been fixed, and no
such claim is needed. The sum of squares imposes all the outside
equations and all the divisions used before taking numerators.

The point substitutions and the pole substitutions have the inverse
maps proved above and below. Among the $212$ fixed chart choices
there is a valid common choice for any list of $211$ source halfpoints.
The $256$ fixed corrections cover the eight Boolean source choices.
A product of branch polynomials permits reuse of the same witness
coordinates, since only one branch must vanish. In particular, no
denominator condition from an inactive branch is imposed.

The coordinate count is
\[
 633+210+3+3+1+1+2+1=854:
\]
point coordinates, pole coordinates, anchor coordinates $(U,v_a,w_a)$,
proxies, $\theta$, $e$, two contact coordinates, and the input.
The final $U$ also enforces the denominators, so no inverse coordinate
is added. All 854 coordinates occur in the supplied specification.

For the bounds, assign each integer constant its absolute coefficient
sum, each variable degree one, and each fixed parameter its maximum
absolute value over its stated finite list. At a sum, degrees take a
maximum and coefficient sums add; at a product, degrees add and
coefficient sums multiply; at a power both bounds scale accordingly.
These exact integer recurrences bound every node. Finally multiply the
branch bounds over $212\cdot256=54272$ branches. Replaying these
recurrences on the supplied circuit gives the displayed bounds.
They are upper bounds, not exact expansion statistics.
\end{proof}

This explicit circuit and the 558-variable indexed construction
serve different purposes. The former can be regenerated directly from
the finite formulas and inspected as a circuit; the latter uses further
finite-algebra elimination whose indexed output is specified and checked,
although its rank-sized matrices and full expansion are not materialized.
Neither is presented as a practical solver for Diophantine equations.
